\scanpage{3}
% archive sheet number: 2
% left half, continuation
\unsure{aussi} que l'ensemble \inserted{ordonné} $\Sigma$ des supports soit ``modulaire'' i.e. $A\wedge\operatorname{Sup}_{i} B_i=\operatorname{Sup}_{i} A\wedge B_i$, donc une ``algèbre de Boole''. Il est possible que cela implique que l'atelier est aussi fidèle.
% slanted
\margin{\figure{\begin{tikzpicture}
\node[rotate=22,align=center] at (0,0) {\unsure{À vérifier}\ ?};
\draw[rotate=22] (0,0) ellipse (1.05cm and .4cm);
\node[anchor=north west] at (-.9,-.35) {\deleted{\unsure{NB}}};
\end{tikzpicture}}}

\uline{Question} Si $x\in\unsure{\mathscr L|\mathcal L}$, alors $x^\circ=\unsure{\operatorname{supp}^\circ|\operatorname{sup}^\circ} x=\operatorname{supp}\unsure{x}\in\Sigma$ est-il un objet minimal de $\Sigma^{\unsure{*}}=\Sigma\setminus\{0_\Sigma\}$ ? En d'autres termes, si $X\in\unsure{\mathscr A|\mathcal A}$, $\deleted{\unsure{X\notin\mathscr L}}$, et si $\operatorname{supp}^\circ X\subset\operatorname{supp}^\circ x$, a-t-on $X=x$. En d'autres termes
\[
X\ne x\overset{?}{\Longrightarrow}\operatorname{supp}^\circ(\unsure{x|X})\unsure{\not\subset}\operatorname{supp}^\circ x
\quad\text{\uuline{i.e.}}\quad
\operatorname{cosupp}^\circ x\unsure{\not\subset}\operatorname{cosupp}^\circ X
\]
i.e. $\exists Z\in\mathscr A$ avec $Z|\circ|x$, $Z\overline{|\circ|}X$ ? Si c'est un local, on se ramène au cas $X=y\in\mathscr L$, (car
\[
X^\circ=\underset{\substack{y\in\mathscr L,\\\text{t.q. }y\unsure{\ll}X\\\text{i.e. }y^\circ\subset X^\circ}}{\operatorname{Sup}}y
\]
), donc à voir si $x,y\in\mathscr L$, $y\ne x$, $\exists Z\in\mathscr A$ avec $Z|\circ|x$, $Z\overline{|\circ|}y$. Je n'y \unsure{ai pensé|arrive} que dans le cas $x\overline{|\circ|}y$.

\margin{\uline{15.7.}}
\uline{Requestion :} Opération de saturation d'un atelier, \unsure{en sens contraire}, faisable dans le cadre des ateliers spatiaux sans plus ? J'en doute !

\diagram{\begin{tikzcd}[ampersand replacement=\&,row sep=small,column sep=small]
\& |[circle,draw,inner sep=1pt]|{S}\arrow[ddl,dash,bend right=30,line width=1.4pt]\arrow[dr,dash]\arrow[ddd,dash]\&\\
\deleted{\unsure{M}}\&\&|[circle,draw,inner sep=1pt]|{LS}\arrow[ddd,dash,line width=1.2pt]\\
|[circle,draw,inner sep=1pt]|{M}\arrow[dddr,dash,bend right=20,line width=1.2pt]\arrow[uur,dash,gray]\&\&\\
\&SF\arrow[dd,dash]\arrow[dr,dash,line width=1.2pt]\&\\
\&\& E\arrow[dd,dash,line width=1pt]\\
\&\unsure{MF}\arrow[dd,dash,line width=1pt]\arrow[dr,dash,line width=1pt]\&\\
\&\&|[circle,draw,inner sep=1pt]|{EM}\\
\&Mq\&
\end{tikzcd}}
\margin{$D\ E\ G\ \unsure{C|E}\ F\ L\ M\ P\ S$\\$Mq,Di,\unsure{v|\nu},p$}

% right half, reused handwritten cover
\hfill\uline{n° 574}

\hfill\uline{300 exemplaires}

Fascicule 0

Recto verso

1) Page de garde

2) Lettre\qquad L 1 à L 43

3) Table des Matières\qquad T 1 : T 10

Les trois dernières pages de la table des matières sont tirées avec le fascicule 4

(à \unsure{600|500} exemplaires)

4) Introduction\qquad (i à cxxii)

\hfill M\textsuperscript{me} LACAN

% slanted lower-left note on the right half
\figure{\begin{tikzpicture}[rotate=28]
\draw (0,0)--(2,0) (0,-.13)--(0,.13) (2,-.13)--(2,.13);
\draw (2,.12)--(4,.12) (2,0)--(2,.25) (4,0)--(4,.25);
\node at (1,-.3){$S$};\node at (3,.43){$T$};
\draw (-.5,-.4)--(-.5,-.7)--(4.4,-.7)--(4.4,-.4);
\end{tikzpicture}}
\margin{$S\wedge T=S\wedge\underbrace{\underbrace{C(1\wedge C^T)}_{C^T}}_{T}=$\\$\deleted{\unsure{A}}$\\$1-C^T$}
